% ==============================================================================
% MÓDULO PRE-CAPÍTULOS: ABSTRACT EN INGLÉS
% Archivo: pre_caps/06_abstract.tex
% Enfoque: 100% Calidad de Atención (Donabedian y Watson)
% ==============================================================================
\section*{ABSTRACT}
\addcontentsline{toc}{section}{ABSTRACT}

\noindent The main objective of this research project is to understand the meaning, structure, and essence of the lived experience of users regarding the quality of healthcare received at the outpatient clinics of Belén Health Center (Category I-3), province of Huamanga, Ayacucho, during the year 2026. Epistemologically grounded in the interpretive-naturalistic paradigm with a phenomenological qualitative design (Edmund Husserl, Max van Manen, and Roberto Hernández-Sampieri), the study is theoretically anchored in Avedis Donabedian's Interpersonal Healthcare Process model and Jean Watson's Human Caring Theory. The research explores the lived experience of caring quality across the four existential lifeworld dimensions (\textit{Lebenswelt}): relationality (warmth of interpersonal encounter, nursing empathy, welcoming hospitality, and communication in Chanka Quechua), temporality (timeliness, unhurried pace, and meaningful time dedicated during consultation), spatiality (visual and auditory privacy, physical comfort, and confidentiality in consultation rooms), and corporeality (respect for bodily modesty and humanized relief of physical suffering and pain). The target population comprises registered outpatient users of the healthcare facility (18,450 inhabitants, 68.2\% under poverty conditions and 72\% active bilingual speakers of Chanka Quechua and Spanish). Informants will be selected through purposeful non-probabilistic qualitative sampling, estimating 12 to 18 key participants governed by the theoretical category saturation principle. Data collection techniques include semi-structured in-depth phenomenological interviews assisted by a validated bilingual interview guide and reflective field observation recorded in research logs. Qualitative data processing will follow Colaizzi's phenomenological method assisted by ATLAS.ti v9 software, identifying essential and divergent categories of healthcare quality. Scientific trustworthiness is established following Guba and Lincoln's criteria (1985): credibility, transferability, dependability, and confirmability, strictly upholding phenomenological reduction (\textit{epoché}) and the bioethical standards of the Declaration of Helsinki.

\vspace{0.8cm}
\noindent \textbf{Keywords (MeSH):} Quality of Health Care; Nursing Care; Interpersonal Relations; Empathy; Humanization of Assistance; Qualitative Research.

\newpage
